\documentclass[12pt,oneside]{report}
\usepackage[utf8]{inputenc}
\usepackage[T1]{fontenc}
\usepackage{lmodern}
\usepackage{microtype}
\usepackage{amsmath,amssymb}
\usepackage{mathtools}
\usepackage{graphicx}
\usepackage{booktabs}
\usepackage{float}
\usepackage{geometry}
\usepackage{hyperref}
\usepackage{setspace}
\setlength{\parindent}{0pt}
\setlength{\parskip}{6pt}
\setstretch{1.15}
\hypersetup{
    colorlinks=true,
    linkcolor=blue,
    urlcolor=blue,
    citecolor=blue
}
\geometry{
    a4paper,
    left=30mm,
    right=25mm,
    top=30mm,
    bottom=30mm,
}

\title{Information and Coding Theory Homework 1}
\author{Yvonne Akinyi Amugaga \\ Neptune: Y2Q07H}
\date{\today}

\begin{document}

\maketitle

\section*{Information and Coding Theory Exercise 1}

We consider an alphabet with five characters:
\[
\{A, B, C, D, E\}
\]
Each text has length 40.

%------------------------------------------------
\section*{1. Nearly Uniform Distribution}

\subsection*{Generated Text (Length 40)}

\texttt{ABCDEBACEDABCDEBACEDABCDEBACEDABCDEBAC}

Each character appears 8 times.

\subsection*{Probability Distribution}

\begin{center}
\begin{tabular}{ccc}
\toprule
Character & Count & Probability \\
\midrule
A & 8 & 0.20 \\
B & 8 & 0.20 \\
C & 8 & 0.20 \\
D & 8 & 0.20 \\
E & 8 & 0.20 \\
\bottomrule
\end{tabular}
\end{center}

\subsection*{Information per Symbol}

\[
I(x) = -\log_2 p(x)
\]

Since $p(x) = 0.2$:

\[
I = -\log_2(0.2)
\]

\[
\log_2(0.2) = \log_2\left(\frac{1}{5}\right) = -\log_2 5
\]

\[
\log_2 5 \approx 2.322
\]

\[
I \approx 2.322 \text{ bits}
\]

\subsection*{Entropy}

\[
H = \sum p(x) \log_2 \frac{1}{p(x)}
\]

\[
H = 5 \times 0.2 \times 2.322
\]

\[
H = 2.322 \text{ bits}
\]

Maximum entropy for 5 symbols:

\[
H_{\max} = \log_2 5 = 2.322
\]

\subsection*{Total Information in the Text}

\[
40 \times 2.322 = 92.88 \text{ bits}
\]

%------------------------------------------------
\section*{2. Biased Distribution (High-Frequency Characters)}

\subsection*{Generated Text (Length 40)}

\texttt{AAAAAAAAAAAAAAAABBBBBBBBCCCCDDDEEE}

Counts:

\[
A = 16,\quad B = 8,\quad C = 6,\quad D = 5,\quad E = 5
\]

\subsection*{Probability Distribution}

\begin{center}
\begin{tabular}{ccc}
\toprule
Character & Count & Probability \\
\midrule
A & 16 & 0.40 \\
B & 8 & 0.20 \\
C & 6 & 0.15 \\
D & 5 & 0.125 \\
E & 5 & 0.125 \\
\bottomrule
\end{tabular}
\end{center}

\subsection*{Information per Symbol}

\[
I(x) = -\log_2 p(x)
\]

\[
I(A) = -\log_2(0.4) \approx 1.322
\]

\[
I(B) = -\log_2(0.2) = 2.322
\]

\[
I(C) = -\log_2(0.15) \approx 2.737
\]

\[
I(D) = I(E) = -\log_2(0.125) = 3
\]

\subsection*{Entropy}

\[
H = \sum p(x) I(x)
\]

\[
H =
0.4(1.322)
+ 0.2(2.322)
+ 0.15(2.737)
+ 0.125(3)
+ 0.125(3)
\]

\[
H =
0.529
+ 0.464
+ 0.411
+ 0.375
+ 0.375
\]

\[
H \approx 2.154 \text{ bits}
\]

\subsection*{Total Information}

\[
40 \times 2.154 = 86.16 \text{ bits}
\]

%------------------------------------------------
\section*{3. Singular and Non-Singular Codes}

\subsection*{Singular Code}

A singular code assigns the same codeword to different symbols.

\begin{center}
\begin{tabular}{cc}
\toprule
Character & Code \\
\midrule
A & 0 \\
B & 0 \\
C & 10 \\
D & 110 \\
E & 111 \\
\bottomrule
\end{tabular}
\end{center}

Here A and B share the codeword 0, therefore the code is singular.

\subsection*{Non-Singular Code}

Each symbol has a distinct codeword.

\begin{center}
\begin{tabular}{cc}
\toprule
Character & Code \\
\midrule
A & 0 \\
B & 10 \\
C & 110 \\
D & 1110 \\
E & 1111 \\
\bottomrule
\end{tabular}
\end{center}

All codewords are distinct, therefore the code is non-singular.
This code is also prefix-free (instantaneous).

%------------------------------------------------
\section*{4. Comparison}

\begin{center}
\begin{tabular}{ccc}
\toprule
Case & Entropy (bits/symbol) & Total Information (40 symbols) \\
\midrule
Uniform & 2.322 & 92.88 \\
Biased & 2.154 & 86.16 \\
\bottomrule
\end{tabular}
\end{center}

Entropy decreases when the distribution becomes less uniform.

\end{document}